\documentclass[../main.tex]{subfiles}
\begin{document}

\section{Problem Statement}
\label{sec:problem-statement}

Understanding human personality is a fundamental goal in psychology and behavioural science, with far-reaching implications for how individuals interact, make decisions, and engage socially. With the proliferation of social media and digital platforms, large volumes of user-generated text have become available --- text that carries latent signals about the author's personality traits. This has created both the opportunity and the need for automated methods that can infer personality from written language.

The applications of personality analysis are wide and consequential. Personalised recommendation systems can adapt content to a user's disposition; mental health platforms can tailor counselling strategies to an individual's emotional profile; human resource tools can assist in team composition; and political analysts can model persuasion susceptibility at scale~\cite{mehta2019recent}. In each case, the quality of the downstream application is bounded by the quality of the personality inference. Automatic personality trait detection from text --- the task of identifying an author's personality traits from their written language using established psychological frameworks --- therefore occupies a strategically important position at the intersection of Natural Language Processing and personality psychology.

The most widely adopted psychological framework for this purpose is the Big Five model, also known as the OCEAN model, which characterises personality along five statistically independent continuous dimensions: Openness to experience, Conscientiousness, Extraversion, Agreeableness, and Neuroticism~\cite{mccrae1992introduction}. These five dimensions have been replicated across languages and cultures and are supported by decades of empirical evidence, making them the natural target for computational personality recognition.

\section{Background and Research Problems}
\label{sec:giaiphap}

Despite sustained research interest spanning two decades, automatic personality recognition from text remains a difficult problem. Early work relied on hand-crafted psycholinguistic features --- most notably the Linguistic Inquiry and Word Count (LIWC) lexicon~\cite{pennebaker2001linguistic} --- and classical classifiers such as Support Vector Machines. These methods demonstrated that writing style correlates reliably with personality, but they cannot capture contextual or compositional meaning. Two texts sharing no vocabulary yet conveying the same psychological content are treated as unrelated; two texts sharing vocabulary yet opposing in meaning are treated as similar.

Deep learning models addressed the representational limitation: convolutional and recurrent networks trained on word embeddings surpassed LIWC-based baselines on standard benchmarks, and fine-tuned transformer models such as BERT pushed accuracy further still. However, the primary evaluation benchmark --- the Essays dataset~\cite{pennebaker1999linguistic} with 2,468 labelled essays --- is small by modern standards. Models trained on it risk capturing dataset-specific stylistic regularities rather than trait-relevant signals, and they generalise poorly across writing genres.

The emergence of large language models (LLMs) such as GPT-4 and LLaMA opened fundamentally different possibilities. Rather than training classifiers from scratch on small labelled corpora, one can leverage the broad world knowledge and reasoning capacity already encoded in a pretrained LLM. Recent work has exploited this in several ways: using LLMs to augment training data through multi-perspective analysis~\cite{hu2024llm}; deploying psychological questionnaire items as chain-of-thought scaffolds~\cite{yang2023psycot}; introducing contrasting personality-induced agents with a judge~\cite{yeo2025pado}; and applying Retrieval-Augmented Generation (RAG) over text embeddings of essays~\cite{ma2025post}.

Each of these approaches advances the state of the art, yet all share a structural limitation: they operate directly on raw text or text-level embeddings, and therefore conflate surface writing style with psychological content. An emotionally expressive introvert and an emotionally expressive extravert produce superficially similar text but carry opposite labels. A student who writes casually about disciplined study habits will appear less conscientious than a student who writes formally about disorganised behaviour, even though the underlying trait signals point in opposite directions. This \emph{surface style confound} is the central structural challenge this thesis addresses. Critically, the goal is not to maximise classification accuracy per se, but to introduce an explicit, interpretable psychological representation that separates surface style from trait-relevant evidence --- a step whose value lies in transparency and diagnostic utility rather than raw F1 improvement.

A related problem affects retrieval-based methods specifically. When training exemplars are retrieved by raw-text similarity, the retrieved neighbours tend to share topical or stylistic features with the query rather than personality-relevant ones. What is needed is retrieval over a representation that is explicitly orthogonal to surface style --- a structured psychological representation. Constructing and evaluating such a representation is the primary technical objective of this work.

\section{Research Objectives and Conceptual Framework}
\label{sec:objectives}

This thesis proposes a three-component pipeline that addresses the surface style confound by interposing an explicit, interpretable psychological representation between raw text and the personality classifier. The primary research question is whether disentangling surface style from structured trait evidence yields a more transparent and diagnosable system --- not whether it yields higher macro-F1 than fine-tuned encoder baselines, which represent a fundamentally different modelling paradigm optimised directly for classification accuracy on this benchmark.

The first component is a \textbf{30-facet psychological profiler}. Rather than classifying text directly, the system first asks an LLM to generate a structured profile of the author using the thirty facets of the NEO-PI-R inventory~\cite{costa2008revised} --- six facets per OCEAN trait. Each facet receives a controlled signal from the vocabulary \{high, mod, low, none, n/e\}, where \textit{n/e} (not evidenced) is a first-class value that prevents the profiler from fabricating signals on facets that the text does not surface. The profile is expressed in a fixed psychological schema rather than in free natural language, which makes it directly comparable across authors and genres.

The second component is a \textbf{hybrid retrieval mechanism}. Each training essay is profiled offline and encoded as a dual embedding that fuses the raw-text embedding with the profile text embedding. At inference time, the test essay is profiled and its top-$k$ nearest training exemplars are retrieved by a hybrid scoring function that combines dense embedding similarity with a trait-masked facet cosine computed over the 30-dimensional ordinal facet vectors. Because the facet cosine measures psychological proximity in the NEO-PI-R space rather than stylistic similarity, retrieved neighbours share personality structure with the query even when their surface style differs.

The third component is a \textbf{retrieval-augmented reasoner}. The retrieved exemplars --- rendered as complete 30-facet profiles with known labels --- serve as calibration anchors for the final LLM classification step. The reasoning prompt enforces a six-block XML structure that tags each piece of evidence as either a stable trait cue or a transient state cue, and tallies only trait cues to produce the final verdict.

\section{Contributions}
\label{sec:contributions}

This thesis makes three contributions. First, it proposes a structured psychological profiling approach that maps raw text to a 30-dimensional NEO-PI-R facet profile, providing a personality-specific intermediate representation that is genre-invariant and directly comparable across authors; the profiler incorporates bias-control mechanisms including an anti-genre-bias clause and the \textit{n/e} signal to suppress fabricated evidence. Second, it introduces a hybrid retrieval method that combines dense dual embeddings with trait-masked ordinal facet cosine similarity, enabling retrieval of psychologically similar training exemplars independent of surface writing style, with per-trait channel weights optimised on a held-out validation set. Third, it presents a structured reasoning prompt (Reasoned-RAG-Def-Oneshot) that uses retrieved 30-facet exemplar profiles as calibration anchors and enforces state-versus-trait cue tagging, and provides a systematic ablation across five retrieval strategies and three pipeline configurations on the Essays benchmark, demonstrating the incremental contribution of each component and identifying state-as-trait leakage on Neuroticism as the primary residual error mode. The ablation results show that each component contributes positively in isolation; the full pipeline does not exceed the macro-F1 of fine-tuned transformer baselines, as the design trades raw discriminative power for interpretability and psychological grounding --- a trade-off examined explicitly in Chapter~\ref{chap:experiments}.

\section{Organisation of Thesis}
\label{sec:organization}

Chapter~\ref{chap:litreview} establishes the theoretical and empirical context. It introduces the Big Five personality model and the NEO-PI-R facet schema, surveys the Essays and MyPersonality benchmark datasets, and traces the evolution of automatic personality recognition from psycholinguistic feature engineering and deep learning through to contemporary LLM-based approaches. The chapter closes by identifying the surface style confound and the absence of structured psychological retrieval as the limitations that motivate this work.

Chapter~\ref{chap:methodology} presents the proposed pipeline in detail. It describes the 30-facet profiler design, the dual embedding construction, the hybrid scoring function and per-trait weight selection, and the structured reasoning prompt.

Chapter~\ref{chap:experiments} reports the experimental evaluation. It covers the experimental setup, dataset statistics, evaluation metrics, a retrieval ablation comparing five strategies on the Essays validation set, a pipeline component ablation, a classifier backbone comparison with prediction bias analysis, and a qualitative error analysis.

\end{document}
